\documentclass[preprint,review,12pt]{elsarticle}
\usepackage[utf8]{inputenc}
\usepackage[T1]{fontenc}
\usepackage{amsmath,amssymb}
\usepackage{graphicx}
\usepackage{booktabs}
\usepackage{textcomp}
\usepackage{siunitx}
\usepackage{url}
\usepackage[hidelinks]{hyperref}
\graphicspath{{figures/}}
\sisetup{group-separator={,}}
% Generated by paper/check_numbers.py --write from results/*.csv. Do not edit by hand.
\newcommand{\mzbTrains}{210}
\newcommand{\mzbTrainKm}{98,602}
\newcommand{\mzbEnergy}{1,292}
\newcommand{\mzbSpecific}{12.8}
\newcommand{\mzbMeanMW}{53.8}
\newcommand{\mzbPeakMW}{151}
\newcommand{\mzbRegen}{68}
\newcommand{\mzbBack}{31}
\newcommand{\mzbBackPct}{2.4}
\newcommand{\mzbIdleMin}{352}
\newcommand{\mzbSubs}{13}
\newcommand{\mzbAccess}{516}
\newcommand{\mzbUsePct}{10}
\newcommand{\mzbAccessMin}{12}
\newcommand{\mzbAccessMax}{66}
\newcommand{\mzbAccessMed}{35}
\newcommand{\mzbTight}{7}
\newcommand{\mzbFirm}{34}
\newcommand{\mzbFirmFlex}{49}
\newcommand{\mzbNoBat}{70}
\newcommand{\mzbFlex}{100}
\newcommand{\mzbBat}{385}
\newcommand{\mzbBatMW}{315}
\newcommand{\mzbBatMWh}{161}
\newcommand{\mzbBatCost}{103}
\newcommand{\mzbFlexCurt}{0.4}
\newcommand{\mzbFirmGrowth}{0}
\newcommand{\mzbNoBatGrowth}{4}
\newcommand{\mzbFlexGrowth}{6}
\newcommand{\mzbBatGrowth}{347}
\newcommand{\mzbAccessTrafo}{1,260}
\newcommand{\mzbFirmTrafo}{595}
\newcommand{\mzbFlexTrafo}{970}
\newcommand{\mzbBatTrafo}{1,056}
\newcommand{\mzbGapS}{59}
\newcommand{\mzbGapMin}{70}
\newcommand{\mzbGapQ}{234}
\newcommand{\mzbNoneSust}{314}
\newcommand{\mzbNoneCut}{71}
\newcommand{\mzbNoneCutPct}{18}
\newcommand{\mzbNoneCutMed}{21}
\newcommand{\mzbBatMinMWh}{116}
\newcommand{\mzbCycMin}{4.1}
\newcommand{\mzbCyc}{2.9}
\newcommand{\mzbRampS}{30}
\newcommand{\mzbRampSMed}{22}
\newcommand{\mzbRampTen}{32}
\newcommand{\mzbPeakS}{45}
\newcommand{\mzbVufMin}{0.2}
\newcommand{\mzbVufMax}{0.8}
\newcommand{\mzbVufWeak}{1.6}
\newcommand{\mzbFinSites}{12}
\newcommand{\mzbFinMW}{379}
\newcommand{\mzbFinIT}{330}
\newcommand{\mzbFinGWh}{2,822}
\newcommand{\mzbCapexDC}{2,966}
\newcommand{\mzbCapexBay}{70}
\newcommand{\mzbCapexBat}{101}
\newcommand{\mzbCapexGen}{146}
\newcommand{\mzbCapexEl}{317}
\newcommand{\mzbGenSaved}{62}
\newcommand{\mzbCapexPV}{52}
\newcommand{\mzbElPct}{11}
\newcommand{\mzbFee}{15.2}
\newcommand{\mzbFeeLow}{7.6}
\newcommand{\mzbFeeHigh}{30.3}
\newcommand{\mzbFeeNPV}{161}
\newcommand{\mzbFeeFirm}{0.8}
\newcommand{\mzbFeeFlex}{3.8}
\newcommand{\mzbBatAnnual}{19.0}
\newcommand{\msTrains}{280}
\newcommand{\msTrainKm}{63,976}
\newcommand{\msEnergy}{723}
\newcommand{\msSpecific}{10.9}
\newcommand{\msMeanMW}{30.1}
\newcommand{\msPeakMW}{90}
\newcommand{\msRegen}{50}
\newcommand{\msBack}{23}
\newcommand{\msBackPct}{3.2}
\newcommand{\msIdleMin}{374}
\newcommand{\msSubs}{12}
\newcommand{\msAccess}{512}
\newcommand{\msUsePct}{6}
\newcommand{\msAccessMin}{35}
\newcommand{\msAccessMax}{53}
\newcommand{\msAccessMed}{42}
\newcommand{\msTight}{1}
\newcommand{\msFirm}{174}
\newcommand{\msFirmFlex}{248}
\newcommand{\msNoBat}{227}
\newcommand{\msFlex}{314}
\newcommand{\msBat}{417}
\newcommand{\msBatMW}{190}
\newcommand{\msBatMWh}{99}
\newcommand{\msBatCost}{63}
\newcommand{\msFlexCurt}{3.6}
\newcommand{\msFirmGrowth}{73}
\newcommand{\msNoBatGrowth}{140}
\newcommand{\msFlexGrowth}{199}
\newcommand{\msBatGrowth}{395}
\newcommand{\msAccessTrafo}{660}
\newcommand{\msFirmTrafo}{292}
\newcommand{\msFlexTrafo}{473}
\newcommand{\msBatTrafo}{549}
\newcommand{\msGapS}{210}
\newcommand{\msGapMin}{227}
\newcommand{\msGapQ}{337}
\newcommand{\msNoneSust}{360}
\newcommand{\msNoneCut}{57}
\newcommand{\msNoneCutPct}{14}
\newcommand{\msNoneCutMed}{10}
\newcommand{\msBatMinMWh}{68}
\newcommand{\msCycMin}{4.7}
\newcommand{\msCyc}{2.8}
\newcommand{\msRampS}{18}
\newcommand{\msRampSMed}{10}
\newcommand{\msRampTen}{19}
\newcommand{\msPeakS}{47}
\newcommand{\msVufMin}{0.2}
\newcommand{\msVufMax}{0.9}
\newcommand{\msVufWeak}{1.9}
\newcommand{\msFinSites}{12}
\newcommand{\msFinMW}{417}
\newcommand{\msFinIT}{363}
\newcommand{\msFinGWh}{3,105}
\newcommand{\msCapexDC}{3,263}
\newcommand{\msCapexBay}{73}
\newcommand{\msCapexBat}{62}
\newcommand{\msCapexGen}{160}
\newcommand{\msCapexEl}{296}
\newcommand{\msGenSaved}{69}
\newcommand{\msCapexPV}{57}
\newcommand{\msElPct}{9}
\newcommand{\msFee}{16.7}
\newcommand{\msFeeLow}{8.3}
\newcommand{\msFeeHigh}{33.4}
\newcommand{\msFeeNPV}{177}
\newcommand{\msFeeFirm}{6.6}
\newcommand{\msFeeFlex}{12.4}
\newcommand{\msBatAnnual}{11.7}
\newcommand{\anAccess}{51}
\newcommand{\anLimit}{45.9}
\newcommand{\anDC}{40}
\newcommand{\anBatMW}{28.8}
\newcommand{\anBatMWh}{14.4}
\newcommand{\anNone}{30.3}
\newcommand{\anCut}{9.7}
\newcommand{\anCutPct}{24}
\newcommand{\anMean}{3.9}
\newcommand{\anUse}{7.6}
\newcommand{\anPeak}{34.7}
\newcommand{\anPeakS}{36.1}
\newcommand{\anGapMin}{11.2}
\newcommand{\anGapQ}{29.6}
\newcommand{\anFirm}{1.5}
\newcommand{\anBridgeMin}{30}
\newcommand{\anFirmBlock}{28}
\newcommand{\allTight}{8}
\newcommand{\resFlatWind}{8.00}
\newcommand{\resFlatPV}{1.50}
\newcommand{\resFlatTot}{9.50}
\newcommand{\resFlatRatio}{2.40}
\newcommand{\resFlatCurt}{0.0}
\newcommand{\resFlatCost}{64}
\newcommand{\resFlexWind}{5.75}
\newcommand{\resFlexPV}{1.00}
\newcommand{\resFlexTot}{6.75}
\newcommand{\resFlexRatio}{1.75}
\newcommand{\resFlexCurt}{2.5}
\newcommand{\resFlexCost}{54}
\newcommand{\resStoWind}{2.75}
\newcommand{\resStoPV}{3.00}
\newcommand{\resStoTot}{5.75}
\newcommand{\resStoRatio}{1.34}
\newcommand{\resStoCurt}{0.0}
\newcommand{\resStoCost}{50}
\newcommand{\resBothWind}{2.50}
\newcommand{\resBothPV}{3.00}
\newcommand{\resBothTot}{5.50}
\newcommand{\resBothRatio}{1.30}
\newcommand{\resBothCurt}{1.7}
\newcommand{\resBothCost}{49}
\newcommand{\resFlatLooseWind}{4.00}
\newcommand{\resFlatLoosePV}{1.50}
\newcommand{\resFlatLooseTot}{5.50}
\newcommand{\resFlatLooseRatio}{1.35}
\newcommand{\resFlatLooseCurt}{0.0}
\newcommand{\resFlatLooseCost}{44}
\newcommand{\resFlexLooseWind}{2.75}
\newcommand{\resFlexLoosePV}{1.25}
\newcommand{\resFlexLooseTot}{4.00}
\newcommand{\resFlexLooseRatio}{1.06}
\newcommand{\resFlexLooseCurt}{8.1}
\newcommand{\resFlexLooseCost}{39}
\newcommand{\resFlexValue}{9.4}
\newcommand{\resFlexValueMzb}{27}

% Generated by paper/check_numbers.py --write from results/*.csv. Do not edit by hand.
\newcommand{\substationrows}{%
\multicolumn{11}{l}{\textit{Madrid--Barcelona}}\\
Villaverde Bajo & 220 & 32 & 2.1 & 32.2 & 36.8 & 0.0 & 0.0 & 0.0 & 26 & 26.0/13.0\\
Anchuelo & 400 & 51 & 3.9 & 34.7 & 44.4 & 1.5 & 11.2 & 16.0 & 40 & 28.8/14.4\\
Brihuega & 400 & 45 & 4.7 & 32.2 & 42.3 & 0.0 & 8.3 & 11.8 & 33 & 24.7/12.4\\
Medinaceli & 400 & 66 & 4.9 & 38.5 & 40.6 & 18.8 & 20.9 & 29.8 & 52 & 31.1/15.6\\
Terrer & 400 & 34 & 4.5 & 39.6 & 41.2 & 0.0 & 0.0 & 0.0 & 24 & 24.0/12.7\\
Rueda de Jal\'on & 400 & 12 & 3.3 & 32.8 & 34.1 & 0.0 & 0.0 & 0.0 & 6 & 6.0/4.2\\
Feria de Muestras & 220 & 35 & 2.1 & 27.8 & 31.5 & 0.0 & 3.7 & 5.3 & 29 & 25.3/14.4\\
Zaragoza & 220 & 34 & 4.2 & 38.3 & 42.0 & 0.0 & 0.0 & 0.0 & 24 & 24.0/12.0\\
Pe\~nalba & 400 & 49 & 5.3 & 27.6 & 35.2 & 8.9 & 16.5 & 23.6 & 36 & 19.5/9.7\\
Lleida & 220 & 35 & 4.8 & 35.1 & 43.9 & 0.0 & 0.0 & 0.0 & 24 & 24.0/12.0\\
L'Espluga & 400 & 55 & 5.2 & 40.3 & 44.2 & 5.3 & 9.2 & 13.2 & 42 & 32.8/16.4\\
La Gornal & 220 & 33 & 5.0 & 35.9 & 48.6 & 0.0 & 0.0 & 0.0 & 23 & 23.0/11.5\\
Sant Boi & 220 & 35 & 3.8 & 39.4 & 53.9 & 0.0 & 0.0 & 0.0 & 26 & 26.0/13.0\\
\textbf{Total} & & \textbf{516} & \textbf{53.8} & & & \textbf{34.5} & \textbf{69.8} & \textbf{99.7} & \textbf{385} & \textbf{315.2/161.3}\\[2pt]
\multicolumn{11}{l}{\textit{Madrid--Seville}}\\
El Hornillo & 220 & 35 & 4.1 & 45.5 & 46.7 & 0.0 & 0.0 & 0.0 & 26 & 26.0/13.0\\
A\~nover & 220 & 43 & 2.0 & 13.7 & 25.2 & 13.5 & 25.0 & 35.7 & 36 & 11.0/5.7\\
Mora & 220 & 45 & 2.5 & 17.7 & 19.2 & 21.3 & 22.8 & 32.6 & 37 & 14.2/7.5\\
El Emperador & 220 & 42 & 2.8 & 15.2 & 20.1 & 17.7 & 22.6 & 32.3 & 34 & 11.4/8.2\\
Ciudad Real & 220 & 45 & 2.9 & 22.3 & 25.3 & 15.2 & 18.2 & 25.9 & 36 & 17.8/8.9\\
La Nava & 220 & 45 & 2.9 & 26.2 & 30.0 & 10.5 & 14.3 & 20.4 & 36 & 21.7/10.9\\
Venta La In\'es & 220 & 53 & 2.8 & 19.5 & 24.5 & 23.1 & 28.1 & 40.1 & 43 & 14.9/7.5\\
Arroyo Valle & 220 & 39 & 2.7 & 22.6 & 26.6 & 8.5 & 12.5 & 17.8 & 31 & 18.5/9.3\\
La Lancha & 220 & 40 & 3.2 & 33.3 & 35.6 & 0.4 & 2.7 & 3.8 & 32 & 29.3/15.2\\
Posadas & 220 & 45 & 1.3 & 7.5 & 21.5 & 19.0 & 33.0 & 40.5 & 38 & 5.0/2.5\\
Lora del R\'{\i}o & 220 & 40 & 1.3 & 8.2 & 11.5 & 24.5 & 27.8 & 36.0 & 34 & 6.2/3.1\\
Rinconada & 220 & 40 & 1.5 & 15.9 & 15.9 & 20.1 & 20.1 & 28.7 & 34 & 13.9/7.0\\
\textbf{Total} & & \textbf{512} & \textbf{30.1} & & & \textbf{173.8} & \textbf{227.1} & \textbf{313.8} & \textbf{417} & \textbf{189.9/98.8}\\[2pt]}

\newcommand{\MEUR}{M\texteuro}

\journal{Applied Energy}

\begin{document}

\begin{frontmatter}

\title{Hosting flexible data centres on the grid-connection capacity of high-speed railway traction substations: hosting capacity, implementation architecture and security for two Spanish lines}

% Author block intentionally left generic in the public repository; complete at submission (see NOTES.md).
\author[a]{Author One\corref{cor}}
\author[b]{Author Two}
\cortext[cor]{Corresponding author.}
\affiliation[a]{organization={Affiliation to be completed at submission}, country={Spain}}
\affiliation[b]{organization={Affiliation to be completed at submission}, country={Spain}}

\begin{abstract}
Grid-connection capacity, not energy, has become the binding constraint for new data centres in Spain, while the traction substations of the high-speed railway network hold transmission-level access permits that they use, on average, at \msUsePct--\mzbUsePct\,\% of their rating. This paper quantifies how much data-centre load could share that capacity with absolute priority for the train, and shows how such an installation can be built and operated safely. A second-by-second traction simulation driven by the published timetables of the three operators is applied to the Madrid--Barcelona (\mzbSubs{} substations, \mzbAccess\,MW granted) and Madrid--Seville (\msSubs{} substations, \msAccess\,MW) lines, using the access capacity actually granted to each railway bay by the transmission system operator. Under current rules only the capacity that survives the loss of a neighbouring substation can be released: \mzbFirm\,MW and \msFirm\,MW. A shared access permit with train priority raises this to \mzbNoBat/\msNoBat\,MW without storage, \mzbFlex/\msFlex\,MW if 30\,\% of the load is curtailable, and \mzbBat/\msBat\,MW with power batteries of \mzbBatMWh/\msBatMWh\,MWh. The averaging window used to enforce the joint limit is the most influential parameter: moving from 1 to 15\,min triples the storage-free capacity on the busier line. We propose an implementation architecture (dedicated bay at the same node, redundant transformers, grid-forming storage, a two-block data centre and a fail-safe limiter), cost it at about 10\,\% of the data-centre investment, and derive physical, cyber and operational security requirements so that the railway acquires no new failure mode. The same flexibility that yields power to the train lowers the cost of hourly renewable matching by about \resFlexValue\,\texteuro/MWh.
\end{abstract}

\begin{highlights}
\item High-speed traction substations use \msUsePct--\mzbUsePct\,\% of their granted grid access on average.
\item A shared access permit with train priority unlocks \mzbBat--\msBat\,MW per line with batteries.
\item Enforcing the joint limit on 15-min instead of 1-min means triples storage-free capacity.
\item A fail-safe, one-way architecture adds no new failure mode to railway operation.
\item Data-centre flexibility cuts the cost of 80\,\% hourly renewable matching by \resFlexValue\,\texteuro/MWh.
\end{highlights}

\begin{keyword}
Data centre \sep Railway traction substation \sep Grid connection capacity \sep Demand flexibility \sep Battery energy storage \sep Critical infrastructure security
\end{keyword}

\end{frontmatter}

\section{Introduction}
\label{sec:intro}

\subsection{Grid access as the scarce resource}

Electricity demand from data centres is growing faster than at any time in the last two decades, driven by the training and serving of artificial-intelligence models \cite{iea2025,masanet2020}. In most mature markets the limiting factor is no longer the price or availability of energy but the right to draw power from the grid at a given node. Spain is a clear example. By 1 October 2026 the transmission system operator (TSO), Red El\'ectrica, had granted 18.5\,GW of demand access on the transmission grid, of which 6.3\,GW corresponded to data centres, none of them yet in service \cite{ree2026mapa}; only about one in four transmission nodes can still accept new demand \cite{ree2026nudos}. The Government has reacted with capacity tenders at saturated nodes and an annual reservation fee \cite{rdl72026}, a new category of flexible access permits \cite{cnmc2026flex}, and a draft royal decree that conditions the access permit of large data centres on efficiency and hourly renewable-matching requirements \cite{prd2026,ashurst2026}.

At the same time, part of the access capacity already granted at transmission level is structurally underused. Railway infrastructure holds 3.55\,GW of granted demand access on the Spanish transmission grid \cite{ree2026mapa}. High-speed traction substations are connected at 220 or 400\,kV and are sized for the worst minutes of the timetable under the loss of an adjacent substation, yet a high-speed line is, electrically, a sparse sequence of short power pulses: the mean load of a substation is a small fraction of its peak. Because the methodology used to compute node margins subtracts the full capacity of every granted permit \cite{cnmc2025capacidad}, this headroom is invisible to other consumers. It is \emph{capacity}, not energy, that is idle: regenerative braking returned to the transmission grid is a minor flow in 25\,kV AC high-speed systems, and this paper does not rely on it.

\subsection{Related work}

Three bodies of literature meet in this problem.

\emph{Data centres as flexible loads.} The potential of data centres for demand response has been studied for more than a decade, from early surveys and pricing mechanisms \cite{wierman2014,liu2014} to spatial load shifting between sites \cite{zheng2020,zhang2020}. Recent work changes the scale of the argument: Norris et al.\ estimate that the United States grid could integrate on the order of 100\,GW of new large loads if these accept curtailment during a fraction of a percent of the year \cite{norris2025}, and a field demonstration on a production GPU cluster sustained a 25\,\% power reduction for three hours without violating quality of service \cite{colangelo2025}. Industry programmes now test this at several sites \cite{epri2024}. Two caveats are equally recent. Training workloads can produce fast, large power swings \cite{li2024}, and voltage-sensitive data centres can disconnect simultaneously: on 10 July 2024 about 1,500\,MW of data-centre load in Virginia transferred to on-site supply after a transmission fault \cite{nerc2025}. Flexibility also interacts with clean-energy procurement: hourly matching is far more demanding than annual matching \cite{xu2024,riepin2024}.

\emph{Electrical behaviour of AC railways.} The traction load is single-phase, pulsating and partly regenerative. Its power-quality impact on the supplying grid, mainly voltage unbalance, is well characterised \cite{gazafrudi2015}, and advanced supply schemes with static converters have been proposed to remove neutral sections and balance the load \cite{krastev2016,pilo2015}. Energy storage and regenerative-braking recovery have been studied extensively, mostly for DC urban systems \cite{gonzalezgil2013,khodaparastan2019,ratniyomchai2014}, together with wider energy-saving measures \cite{douglas2015}. Stationary storage and photovoltaic generation have been co-optimised with the loads of a railway station \cite{sengor2018}.

\emph{Sharing railway assets with the power system.} Rail corridors have been proposed as carriers of mobile storage \cite{moraski2023}. Closest to this work, a recent study estimates the non-railway load that the DC traction substations of the Rotterdam metro could host, between 2 and 5\,MW, and finds that the binding limit is sometimes the contractual grid-connection capacity rather than the equipment \cite{dwarkasing2026}. In Spain, the competition authority has recommended that the infrastructure manager publish the location and consumption of its high-speed substations and offer adjacent land for renewable generation \cite{cnmc2023adif}.

\subsection{Gap and contributions}

To the best of our knowledge, no study has quantified the data-centre load that high-speed (25\,kV AC, transmission-connected) traction substations could host by sharing their grid-access capacity, nor specified how such an installation should be built, protected and regulated. The question has become more demanding in 2026. Two serious accidents on the Madrid--Seville high-speed line, in January and October \cite{adamuz2026,urda2026}, repeated cable thefts \cite{cable2026} and a large data breach at the infrastructure manager and the incumbent operator \cite{cyber2026} have put railway safety and security at the centre of public debate and led to changes at the head of the transport ministry and of the infrastructure manager \cite{infobae2026}. None of these events involved the traction power supply, but they set the standard that any proposal for third-party use of railway electrical infrastructure must meet: it must be demonstrably unable to degrade railway operation.

The contributions of this paper are:
\begin{enumerate}
\item A reproducible, open model that derives second-by-second substation loads from public timetables and maps them onto the access capacity actually granted to each railway bay, for a high-density line (Madrid--Barcelona) and a low-density line (Madrid--Seville).
\item The hosting capacity for data centres under current rules and under a proposed \emph{shared access permit with priority}, without storage, with a curtailable share and with optimally sized batteries, including sensitivity to traffic growth, access level and the averaging window of the joint limit.
\item An implementation architecture, with its control hierarchy, cost and a staged pilot, showing that the concept can be built with standard equipment.
\item A security analysis (train priority, ride-through, systemic risk, physical and cyber security) that turns the lessons of recent railway and grid incidents into design requirements.
\item The regulatory pathway and the interaction with the Spanish data-centre decree, including the cost of hourly renewable matching with and without flexibility.
\end{enumerate}

Section~\ref{sec:methods} describes the system and the model; Section~\ref{sec:results} gives the results; Section~\ref{sec:arch} presents the architecture and its feasibility; Section~\ref{sec:security} deals with security and resilience; Section~\ref{sec:discussion} discusses regulation, limitations and future work; Section~\ref{sec:conclusions} concludes.

\section{System description and methodology}
\label{sec:methods}

\subsection{System under study}

Spanish high-speed lines are electrified at 25\,kV, 50\,Hz. The Madrid--Zaragoza--Barcelona line (621\,km) uses the $2\times25$\,kV autotransformer system; the Madrid--C\'ordoba--Seville line (471\,km), opened in 1992, uses $1\times25$\,kV. Each traction substation takes two phases of the three-phase transmission grid through a dedicated bay in a TSO substation at 220 or 400\,kV and feeds two catenary sections separated by neutral zones. Typical ratings are two 60\,MVA transformers per substation on the first line and two 30\,MVA units on the second; the latter are being renewed. The infrastructure manager is, legally, an ordinary transmission-connected consumer that holds one access permit per bay \cite{convenio2021}.

Table~\ref{tab:substations} lists the \mzbSubs{} substations of the Madrid--Barcelona line and the \msSubs{} of the Madrid--Seville line with the access capacity granted to each railway bay, taken from the list of occupied demand capacity per bay published by the TSO on 1 October 2026 \cite{ree2026posicion}. Granted capacity ranges from \mzbAccessMin{} to \mzbAccessMax\,MW on the first line (median \mzbAccessMed\,MW) and from \msAccessMin{} to \msAccessMax\,MW on the second. For six substations no bay could be identified in the TSO list (they are probably fed from the distribution grid) and a value of 35--40\,MW, close to the line median, is assumed. Substation chainages are not public and were estimated from geography with an uncertainty of about $\pm10$\,km.

Two lines were chosen to span the range of traffic density. Madrid--Barcelona is the busiest corridor of the network, served by three operators. Madrid--Seville carries less end-to-end traffic but its first kilometres are shared with trains to Toledo, the Levante corridor and M\'alaga.

\subsection{Overview of the model}

The model is a chain of five steps, each a script in the repository: (i) timetable extraction, (ii) traction simulation, (iii) hosting-capacity analysis, (iv) architecture and contingency indicators and cost model, and (v) hourly renewable matching. All parameters that are not verified public data are flagged as assumptions (Table~\ref{tab:assumptions}).

\subsection{Timetable}

The simulated day is Wednesday 21 October 2026. Renfe services are read from the operator's open GTFS feed (downloaded on 12 October 2026). Every trip that uses any part of the line is kept, including trains that join or leave it at junctions. The other two operators are not in that feed: Iryo services follow their published weekday timetable of October 2026 (14 trains per direction Madrid--Barcelona; 7 and 6 Madrid--Seville; 4 per direction Madrid--M\'alaga), and Ouigo services follow the announced frequencies (6, 4 and 3 per direction, respectively) with departures spread evenly between 07:00 and 20:00, which is an assumption. The resulting day has \mzbTrains{} trains and \mzbTrainKm{} train-km on Madrid--Barcelona and \msTrains{} trains and \msTrainKm{} train-km on Madrid--Seville.

\subsection{Traction simulation}

Each train is simulated at 1\,s resolution between consecutive timetable stops. The speed profile consists of acceleration at maximum tractive effort, cruising at the speed $v_c$ that meets the scheduled run time reduced by a regularity margin of 8\,\%, and service braking at 0.5\,m/s$^2$. The power at the wheel is
\begin{equation}
P_w(t) = \left[\,\rho\, m\, a(t) + A + B\,v(t) + C\,v(t)^2 + m\,g\,s(x)\,\right] v(t),
\label{eq:wheel}
\end{equation}
where $m$ is the laden mass, $\rho=1.06$ the rotating-mass factor, $a$ the acceleration, $A$, $B$, $C$ the Davis running-resistance coefficients \cite{rochard2000} and $s(x)$ the gradient at position $x$, obtained from an approximate elevation profile ($\pm50$\,m). $P_w$ is limited to the rated power of the train in traction and in electrical braking. The power at the pantograph is
\begin{equation}
P_p(t) =
\begin{cases}
P_w(t)/\eta_t + P_{aux}, & P_w(t) \ge 0,\\
\eta_r\,P_w(t) + P_{aux}, & P_w(t) < 0,
\end{cases}
\label{eq:panto}
\end{equation}
with $\eta_t=0.85$, $\eta_r=0.80$ and auxiliary power $P_{aux}$ of 0.2--0.5\,MW depending on the train type. Five rolling-stock classes are parameterised from public data sheets (Table~\ref{tab:assumptions}).

Each substation $k$ supplies the trains located between the midpoints to its neighbours, where the neutral sections are assumed to be:
\begin{equation}
P_k(t) = \sum_{i:\,x_i(t)\in Z_k} P_{p,i}(t).
\label{eq:sub}
\end{equation}
Regenerated power is first used by other trains in the same section; the remainder flows back to the grid. Transmission losses in the catenary are not modelled. The load used in the analysis is $L_k(t)=\max(P_k(t),0)$, averaged over windows of 1\,s, 1\,min or 15\,min as needed.

As a plausibility check, the simulated specific consumption is \mzbSpecific\,kWh per train-km on Madrid--Barcelona and \msSpecific\,kWh per train-km on Madrid--Seville, in the range reported for high-speed operation \cite{douglas2015}; the energy returned to the transmission grid is \mzbBackPct\,\% and \msBackPct\,\% of the energy drawn, which confirms that regeneration is a second-order effect at the grid connection.

\subsection{Hosting capacity}

Let $C_k$ be the access capacity of substation $k$ and $\sigma=0.9$ a safety factor. The data centre connects to the same transmission node through its own bay and transformers (Section~\ref{sec:arch}), so the only shared quantity is the joint power limit. The headroom available to the data centre is
\begin{equation}
H_k(t) = \sigma\,C_k - L_k(t).
\label{eq:headroom}
\end{equation}
If a neighbouring substation is out of service, substation $k$ is assumed to take half of the load of its most loaded neighbour:
\begin{equation}
L^{N-1}_k(t) = L_k(t) + \tfrac{1}{2}\max\left(L_{k-1}(t),\,L_{k+1}(t)\right).
\label{eq:n1}
\end{equation}
Five options are evaluated, all with 1-min averages:
\begin{align}
D^{firm}_k &= \max\left(0,\ \sigma C_k - \max_t L^{N-1}_k(t)\right), \label{eq:firm}\\
D^{firm+flex}_k &= \min\left(D^{firm}_k/(1-\phi),\ \sigma C_k\right), \label{eq:firmflex}\\
D^{shared}_k &= \min_t H_k(t), \label{eq:shared}\\
D^{flex}_k &= \min\left(D^{shared}_k/(1-\phi),\ \sigma C_k\right), \label{eq:flex}
\end{align}
where $\phi=0.3$ is the share of the data-centre load that can be curtailed on request (training and batch workloads). $D^{firm}$ is the capacity that the railway never uses, even under N-1; it is the only quantity compatible with current rules, through a voluntary reduction of the railway permit (Section~\ref{sec:regulation}). The other options require the joint limit to be shared in real time.

The fifth option adds a battery that supplies the data centre whenever $H_k(t)<D$. For a constant load $D$, the required battery power and energy are
\begin{align}
P^{b}_k(D) &= \max_t\left(D - H_k(t)\right)^+, \label{eq:pb}\\
E^{b}_k(D) &= \max_t\left[\max_{\tau\le t} S(\tau) - S(t)\right],\qquad S(t)=\sum_{u\le t}\left(H_k(u)-D\right)\Delta t, \label{eq:eb}
\end{align}
that is, the largest cumulative deficit of the day treated as cyclic, which follows from the Lindley recursion of a single-server queue \cite{lindley1952}. Round-trip losses are neglected. To limit degradation the energy is constrained to $E^b_k \ge P^b_k/2$ (a maximum rate of 2C). The size $D^{bat}_k$ is increased in steps of 1\,MW while the marginal battery cost remains below $v=1$\,\MEUR{} per additional MW of data-centre capacity, with battery costs of 0.20\,\MEUR/MW plus 0.25\,\MEUR/MWh, of the order of published projections \cite{schmidt2019}.

Access capacity is varied in five scenarios: the capacity granted today (base case), the rating of the traction transformers (about 100\,MW with $2\times60$\,MVA and 55\,MW with $2\times30$\,MVA), and uniform values of 30, 60 and 120\,MW. Traffic growth and consumption uncertainty are represented by scaling the load by 1.5.

\subsection{Architecture and contingency indicators}

For the battery option the model computes, per substation: the equivalent full cycles per day; the capacity that the same battery sustains under Eq.~\eqref{eq:n1}, hence the curtailment required when a neighbour trips; the storage-free capacity of Eq.~\eqref{eq:shared} for averaging windows of 1\,s, 1\,min and 15\,min; the largest 1\,s and 10\,s changes of the traction load; and an upper bound of the voltage unbalance factor at the transmission node, $\mathrm{VUF}\approx \hat{L}_k/S_{sc}$, with the 1\,s peak $\hat{L}_k$ and an assumed short-circuit power $S_{sc}$ of 12\,GVA at 400\,kV and 5\,GVA at 220\,kV (half of these values for a weak node).

\subsection{Cost model}

Sites below 10\,MW are discarded. Electrical capital cost comprises the dedicated high-voltage bay (3\,\MEUR{} per site), two transformers each rated for the full load (0.04\,\MEUR/MVA), the battery, and medium-voltage generating sets with $N{+}1$ redundancy for the firm block only (0.5\,\MEUR/MW). The data centre itself is costed at 9\,\MEUR{} per MW of IT load with a power usage effectiveness (PUE) of 1.15. The infrastructure manager is assumed to receive an access-sharing fee of 40\,\texteuro/kW per year (range 20--80), discounted at 7\,\% over 20 years. All unit costs are order-of-magnitude assumptions and are listed, with ranges, in the repository.

\subsection{Hourly renewable matching}

The draft decree requires that, in every hour, at least 80\,\% of consumption be matched by new renewable generation \cite{prd2026}. For a data centre of 1\,MW we search the wind and solar portfolio of least cost such that the share of hours failing the 80\,\% test does not exceed 5\,\% (lowest surcharge tier) or 20\,\% (permit-loss threshold). Hourly shapes are those of peninsular wind and solar generation in 2025 \cite{ree2025horario}, scaled to capacity factors of 0.26 and 0.21. Four cases are compared: a flat load; a load that may reduce up to 30\,\% in deficit hours; a 4\,h shifting battery; and both. Surplus is valued at 30\,\texteuro/MWh.

\begin{table}[t]
\caption{Main parameters and assumptions. A: assumption; P: public data.}
\label{tab:assumptions}
\centering\footnotesize
\begin{tabular}{p{0.34\linewidth}p{0.36\linewidth}p{0.2\linewidth}}
\toprule
Parameter & Value & Status\\
\midrule
Simulated day & Wednesday 21 October 2026 & P (timetables)\\
Ouigo departure times & evenly spread 07:00--20:00 & A\\
Rolling stock (mass / power) & AVE 425\,t / 8.8\,MW; Iryo 500\,t / 9.8\,MW; Ouigo 430\,t / 9.3\,MW; Alvia 320\,t / 4.8\,MW; Avant 290\,t / 4.4\,MW & P (data sheets), laden mass A\\
Efficiencies $\eta_t$ / $\eta_r$ & 0.85 / 0.80 & A\\
Regularity margin & 8\,\% & A\\
Elevation profile & approximate, $\pm50$\,m & A\\
Substation chainage & estimated, $\pm10$\,km & A\\
Granted access $C_k$ & TSO list of 1 Oct 2026; 35--40\,MW where unknown (6 substations) & P / A\\
Safety factor $\sigma$ & 0.9 & A\\
N-1 load transfer & 50\,\% of the most loaded neighbour & A\\
Curtailable share $\phi$ & 30\,\% & A\\
Battery cost; C-rate & 0.20\,\MEUR/MW + 0.25\,\MEUR/MWh; $\le 2$C & A\\
Marginal value $v$ & 1\,\MEUR{} per MW & A\\
Short-circuit power & 12\,GVA (400\,kV), 5\,GVA (220\,kV) & A\\
PUE & 1.15 & draft decree\\
\bottomrule
\end{tabular}
\end{table}

\section{Results}
\label{sec:results}

\subsection{Traction load}

On the simulated day the Madrid--Barcelona line draws \mzbEnergy\,MWh from the transmission grid, with a mean of \mzbMeanMW\,MW and a 1-min peak of \mzbPeakMW\,MW for the whole line; Madrid--Seville draws \msEnergy\,MWh, with a mean of \msMeanMW\,MW and a peak of \msPeakMW\,MW (Fig.~\ref{fig:profile}a). Against the \mzbAccess\,MW and \msAccess\,MW of granted access, the mean utilisation is \mzbUsePct\,\% and \msUsePct\,\%. Both lines are electrically idle for about six hours every night (\mzbIdleMin{} and \msIdleMin{} minutes with less than 1\,MW), which coincides with the maintenance window.

The picture at substation level is different (Fig.~\ref{fig:profile}b and Fig.~\ref{fig:substations}). The load is a sequence of pulses of a few minutes: at Anchuelo, the first 400\,kV substation east of Madrid, the mean is \anMean\,MW (\anUse\,\% of the \anAccess\,MW granted) but the 1-min peak reaches \anPeak\,MW. In \allTight{} of the 25 substations the simulated 1-min peak exceeds 90\,\% of the granted access, \mzbTight{} of them on Madrid--Barcelona. This indicates that the infrastructure manager has requested access close to its real peaks, or that the estimated chainages assign too much load to some substations; either way, the firm margin inside the current permits is small on the busier line. The traction load changes by up to \mzbRampS\,MW in one second (median of the per-substation maxima: \mzbRampSMed\,MW on Madrid--Barcelona and \msRampSMed\,MW on Madrid--Seville), which determines the speed required from any device that enforces the joint limit.

\begin{figure}[t]
\centering
\includegraphics[width=\linewidth]{fig_profile}
\caption{Simulated traction demand on 21 October 2026: (a) both lines; (b) one substation against its granted access capacity and the joint limit.}
\label{fig:profile}
\end{figure}

\begin{figure}[t]
\centering
\includegraphics[width=\linewidth]{fig_substations}
\caption{Granted access capacity, simulated 1-min peaks (normal operation and with a neighbouring substation out of service) and daily mean per substation.}
\label{fig:substations}
\end{figure}

\subsection{Hosting capacity}

Table~\ref{tab:substations} and Fig.~\ref{fig:capacity} give the capacity available to data centres for the five options.

\emph{Under current rules} only the firm capacity can be released: \mzbFirm\,MW on Madrid--Barcelona, concentrated in four substations, and \msFirm\,MW on Madrid--Seville, where the lighter traffic leaves margin almost everywhere. If the data centre accepts curtailment of 30\,\% of its load under N-1, these values rise to \mzbFirmFlex\,MW and \msFirmFlex\,MW.

\emph{With a shared access permit}, the data centre can use the headroom of normal operation and yield when a contingency occurs. Without storage the capacity is \mzbNoBat\,MW and \msNoBat\,MW; with a 30\,\% curtailable share, \mzbFlex\,MW and \msFlex\,MW. Curtailment is rare: the flexible share is called on for at most \mzbFlexCurt\,\% and \msFlexCurt\,\% of its energy in the worst substation of each line.

\emph{With batteries}, the capacity reaches \mzbBat\,MW on Madrid--Barcelona (\mzbBatMW\,MW / \mzbBatMWh\,MWh of storage, \mzbBatCost\,\MEUR) and \msBat\,MW on Madrid--Seville (\msBatMW\,MW / \msBatMWh\,MWh, \msBatCost\,\MEUR). The two lines converge to about 75--80\,\% of the granted access, but by different routes: the busy line needs a battery almost as large as the data centre in power terms, because trains fill the permit for a few minutes many times a day, whereas the quiet line needs less than half.

\begin{figure}[t]
\centering
\includegraphics[width=0.8\linewidth]{fig_capacity}
\caption{Data-centre capacity that can share the granted access of the traction substations, by option. Plain bars: Madrid--Barcelona; hatched bars: Madrid--Seville.}
\label{fig:capacity}
\end{figure}

\begin{table}[p]
\caption{Substations, granted access and hosting capacity (base case: granted access, traffic $\times1.0$). Powers in MW; battery in MW/MWh. Mean and peaks are simulated 1-min traction loads.}
\label{tab:substations}
\centering\scriptsize
\begin{tabular}{lrrrrrrrrrr}
\toprule
Substation & kV & $C_k$ & Mean & Peak & Peak N-1 & Firm & Shared & Sh.+flex & Sh.+bat. & Battery\\
\midrule
\substationrows
\bottomrule
\end{tabular}
\end{table}

\subsection{Storage is a power application}

The battery that shaves the train peaks is very different from the energy-shifting batteries used with renewables. If it were sized with the minimum energy of Eq.~\eqref{eq:eb} (\mzbBatMinMWh\,MWh and \msBatMinMWh\,MWh), it would discharge fully in 15--20\,min and perform a median of \mzbCycMin{} and \msCycMin{} equivalent full cycles per day, which would exhaust a lithium iron phosphate system rated for several thousand cycles \cite{preger2020} in three to four years. With the 2C constraint the installed energy rises to \mzbBatMWh\,MWh and \msBatMWh\,MWh and the median cycling falls to \mzbCyc{} and \msCyc{} cycles per day (Fig.~\ref{fig:battery}), consistent with a service life of six to eight years; the cost model uses eight. Chemistries with higher cycle life, or a hybrid with supercapacitors or flywheels for the first seconds \cite{ratniyomchai2014}, are natural alternatives.

\begin{figure}[t]
\centering
\includegraphics[width=\linewidth]{fig_battery}
\caption{(a) Battery energy required as a function of the data-centre size at each substation (grey curves) and selected sizes. (b) Installed energy and median daily cycling with the minimum energy and with the 2C limit.}
\label{fig:battery}
\end{figure}

\subsection{The averaging window of the joint limit}

How the joint limit is measured matters more than any other parameter (Fig.~\ref{fig:window}). Without storage, the capacity on Madrid--Barcelona is \mzbGapS\,MW if the limit is enforced on 1\,s values, \mzbGapMin\,MW on 1-min means and \mzbGapQ\,MW on 15-min means; on Madrid--Seville, \msGapS, \msGapMin{} and \msGapQ\,MW. Contracted power is billed on quarter-hourly maxima, and the capacity methodology itself admits overloads of 15\,\% for up to 20\,min after a contingency \cite{cnmc2025capacidad}; equipment protection, in contrast, acts in seconds. A regulatory definition of the shared permit based on 15-min means, with a fast backstop limiter set at the thermal capability of the bay, would therefore deliver on the busy line more than three times the storage-free capacity of a 1-min rule.

\begin{figure}[t]
\centering
\includegraphics[width=0.6\linewidth]{fig_window}
\caption{Storage-free hosting capacity as a function of the averaging window used to enforce the joint limit.}
\label{fig:window}
\end{figure}

\subsection{Sensitivity to access level and traffic growth}

Fig.~\ref{fig:sensitivity} shows the firm, flexible and battery options for the five access scenarios and for a 50\,\% higher traction load. Three results stand out. First, with the access granted today the firm capacity on Madrid--Barcelona disappears if traffic grows by half (\mzbFirmGrowth\,MW; \msFirmGrowth\,MW remain on Madrid--Seville), whereas the battery option is robust (\mzbBatGrowth\,MW and \msBatGrowth\,MW): sharing with storage is the only option that survives traffic growth without new access. Second, if access were extended to the rating of the existing traction transformers (\mzbAccessTrafo\,MW and \msAccessTrafo\,MW in total), firm capacity would rise to \mzbFirmTrafo\,MW and \msFirmTrafo\,MW and the battery option to \mzbBatTrafo\,MW and \msBatTrafo\,MW; this expansion depends on the margin of each transmission node and is outside the scope of the shared permit. Third, the uniform 30\,MW scenario leaves nothing without storage on the busy line, which confirms that the granted values are close to the real need.

\begin{figure}[t]
\centering
\includegraphics[width=\linewidth]{fig_sensitivity}
\caption{Sensitivity of the hosting capacity to the access capacity per substation and to traffic growth (dashed: traction load $\times1.5$).}
\label{fig:sensitivity}
\end{figure}

\subsection{Hourly renewable matching}

A flat data-centre load needs \resFlatWind\,MW of new wind and \resFlatPV\,MW of new solar capacity per MW (\resFlatTot\,MW in total, \resFlatRatio{} times its consumption in energy) to meet the 80\,\% hourly test in 95\,\% of the hours (Fig.~\ref{fig:renewable}), at a net cost of \resFlatCost\,\texteuro/MWh. If the flexible block reduces consumption in deficit hours, \resFlexTot\,MW suffice and the net cost falls to \resFlexCost\,\texteuro/MWh, with \resFlexCurt\,\% of the energy curtailed. A 4\,h shifting battery brings the portfolio to \resStoTot\,MW (\resStoCost\,\texteuro/MWh), and both together to \resBothTot\,MW (\resBothCost\,\texteuro/MWh). The flexibility that the data centre needs anyway to yield power to the train is thus worth about \resFlexValue\,\texteuro/MWh in compliance cost, or some \resFlexValueMzb\,\MEUR{} per year for the \mzbFinGWh\,GWh consumed by the Madrid--Barcelona fleet. These figures are optimistic, because the aggregate peninsular profile is smoother than that of any real portfolio, but the ranking of the cases is robust and consistent with system-level studies of hourly matching \cite{xu2024,riepin2024}.

\begin{figure}[t]
\centering
\includegraphics[width=0.6\linewidth]{fig_renewable}
\caption{New renewable capacity needed per MW of data centre to meet 80\,\% hourly matching in 95\,\% of the hours, and net cost per MWh consumed.}
\label{fig:renewable}
\end{figure}

\section{Implementation architecture and feasibility}
\label{sec:arch}

The results above are only useful if the installation can be built with standard equipment, under an arrangement in which the railway keeps full control of its own supply. This section specifies such an installation.

\subsection{Connection scheme}

Fig.~\ref{fig:sld} shows the proposed single-line diagram. Its defining choice is that the data centre does \emph{not} connect behind the railway meter or to the railway busbar. It has its own bay in the same TSO substation, its own step-down transformers and its own medium-voltage network. The only thing shared is the joint power limit at the node. Table~\ref{tab:design} summarises the design decisions.

\begin{figure}[t]
\centering
\includegraphics[width=\linewidth]{fig_sld}
\caption{Proposed single-line diagram. The railway bay and the data-centre bay are electrically independent; they share only the joint power limit, enforced by a limiter that reads the power of the railway bay and can trip data-centre feeders.}
\label{fig:sld}
\end{figure}

\begin{table}[t]
\caption{Design decisions of the proposed installation.}
\label{tab:design}
\centering\footnotesize
\begin{tabular}{p{0.16\linewidth}p{0.26\linewidth}p{0.2\linewidth}p{0.28\linewidth}}
\toprule
Element & Choice & Rejected alternative & Reason\\
\midrule
Point of connection & Dedicated bay at the same transmission node & Tap on the railway busbar or line & Required today for flexible access; no common-mode failure with the railway supply\\
Transformers & $2\times100\,\%$, 400(220)/30\,kV & One unit or $2\times50\,\%$ & N-1 and concurrent maintenance (Tier III)\\
Distribution & 30\,kV ring & 20\,kV & Lower currents for tens of MW\\
Storage & Medium-voltage, AC-coupled, grid-forming & Inside the UPS DC link & Serves train-priority shaving and bridges supply losses; can hold the 30\,kV island\\
Photovoltaics & Optional, DC-coupled to the battery, no export & Plant with grid export & Adds energy, not capacity; no generation access permit needed\\
Generating sets & Medium voltage, $N{+}1$, firm block only & Running them to cover train-priority curtailment & Emission-limited hours; would amount to unlicensed generation behind the limit\\
\bottomrule
\end{tabular}
\end{table}

All components are catalogue equipment. Dedicated transmission bays and 400/30\,kV transformers of 50--100\,MVA are standard for wind and solar plants; medium-voltage grid-forming battery systems are an established technology \cite{lin2020}; and medium-voltage generating sets are the norm in large data centres.

\subsection{Two-block data centre}

The data centre is divided into two blocks.

The \emph{firm block} (about 70\,\% of the load: inference, storage, colocation) is a conventional concurrently maintainable design with UPS and generating sets. Its certification is unaffected by the shared connection, because in the Tier classification the utility supply is regarded as an economic alternative and the on-site engine generators as the reliable source \cite{uptime2018}.

The \emph{flexible block} (about 30\,\%: model training, batch and high-performance computing) has no generating sets. It reduces power by capping accelerators within seconds and by checkpointing and pausing jobs within minutes, as demonstrated on production clusters \cite{colangelo2025}. Not installing generators for this block saves \mzbGenSaved\,\MEUR{} and \msGenSaved\,\MEUR{} on the two lines, more than half of the cost of the battery.

\subsection{Control hierarchy}

Train priority is enforced by four layers that act on different time scales:
\begin{enumerate}
\item \emph{Milliseconds to one second: battery.} The grid-forming battery supplies the difference whenever the measured railway power plus the data-centre power approaches the limit. It is the only resource fast enough for load steps of \mzbRampS\,MW in one second, and it also filters the power oscillations of training workloads \cite{li2024}.
\item \emph{One to ten seconds: power capping} of the accelerators in the flexible block.
\item \emph{Minutes: workload management}, pausing or migrating flexible jobs, including to other sites along the line \cite{zheng2020}.
\item \emph{Hard backstop: limiter relay.} An independent protection relay trips feeders of the flexible block, and ultimately of the whole data centre, if the joint power exceeds the limit for longer than a set time. This relay, not the energy management system, is the contractual guarantee of train priority (Section~\ref{sec:security}).
\end{enumerate}
The only input from the railway side is the active power of the railway bay, which the TSO already measures at the bay; no signal flows from the data centre to the railway.

Fig.~\ref{fig:timeline} illustrates the sequence at Anchuelo when a neighbouring substation trips. The site hosts \anDC\,MW with a battery of \anBatMW\,MW / \anBatMWh\,MWh. With the neighbour out, the same battery sustains \anNone\,MW, so \anCut\,MW (\anCutPct\,\%) must be shed: the battery covers the step instantly, power capping follows within seconds and job pausing within minutes. The firm block (\anFirmBlock\,MW) is untouched.

\begin{figure}[t]
\centering
\includegraphics[width=0.6\linewidth]{fig_timeline}
\caption{Illustrative response at Anchuelo after the loss of a neighbouring traction substation. Power levels are model results; the timing of the capping and pausing steps is indicative.}
\label{fig:timeline}
\end{figure}

\subsection{Power quality}

The traction load is single-phase and the data centre is a balanced three-phase load, so the latter does not worsen the unbalance, but it is exposed to it. With the assumed short-circuit powers the upper bound of the voltage unbalance factor at the transmission node is \mzbVufMin--\mzbVufMax\,\% on Madrid--Barcelona and \msVufMin--\msVufMax\,\% on Madrid--Seville, and up to \msVufWeak\,\% in a weak 220\,kV node. Typical planning levels at high voltage are of the order of 1.4\,\% \cite{iec61000313}; the bound is conservative because the V-connection of the traction transformers and the phase rotation between consecutive substations reduce the net unbalance \cite{gazafrudi2015}. Active-front-end UPS rectifiers tolerate these levels, but a joint power-quality study with the real short-circuit power of each node, including harmonics and resonances, is required per site.

\subsection{Cost}

Table~\ref{tab:cost} gives the electrical investment for the battery option. The complete electrical installation costs \mzbCapexEl\,\MEUR{} on Madrid--Barcelona and \msCapexEl\,\MEUR{} on Madrid--Seville, about \mzbElPct\,\% and \msElPct\,\% of the investment in the data centres themselves. For the infrastructure manager, a fee of 40\,\texteuro/kW per year would yield \mzbFee\,\MEUR{} and \msFee\,\MEUR{} per year (\mzbFeeLow--\mzbFeeHigh{} and \msFeeLow--\msFeeHigh\,\MEUR{} over the assumed range), with a present value of \mzbFeeNPV\,\MEUR{} and \msFeeNPV\,\MEUR{}; under current rules the firm capacity alone would yield \mzbFeeFirm\,\MEUR{} and \msFeeFirm\,\MEUR{} per year. These revenues require no railway investment and no change in railway operation.

\begin{table}[t]
\caption{Electrical investment for the option with shared access and batteries (sites of 10\,MW or more).}
\label{tab:cost}
\centering\footnotesize
\begin{tabular}{lrr}
\toprule
 & Madrid--Barcelona & Madrid--Seville\\
\midrule
Sites / data-centre power (MW) & \mzbFinSites{} / \mzbFinMW & \msFinSites{} / \msFinMW\\
Dedicated bays and transformers (\MEUR) & \mzbCapexBay & \msCapexBay\\
Batteries, 2C (\MEUR) & \mzbCapexBat & \msCapexBat\\
Generating sets, firm block, $N{+}1$ (\MEUR) & \mzbCapexGen & \msCapexGen\\
\textbf{Total electrical (\MEUR)} & \textbf{\mzbCapexEl} & \textbf{\msCapexEl}\\
\midrule
Saving: no generators for the flexible block (\MEUR) & \mzbGenSaved & \msGenSaved\\
Optional photovoltaics, 0.25\,MWp per MW (\MEUR) & \mzbCapexPV & \msCapexPV\\
For reference: data centres (\MEUR) & \mzbCapexDC & \msCapexDC\\
\bottomrule
\end{tabular}
\end{table}

\subsection{Staged pilot}

The concept can be validated at a single site before any regulatory change, in three stages.

\emph{Stage 1: measurement (no construction).} One year of 1\,s power records of a railway bay, compared with the simulation, would settle the largest uncertainty of this study. Anchuelo is a suitable candidate: it is a 400\,kV node close to Madrid with \anAccess\,MW granted, a simulated mean of \anMean\,MW and a firm margin of \anFirm\,MW.

\emph{Stage 2: firm and flexible connection under current rules.} A data centre of the order of 10\,MW with its own bay and a flexible access permit \cite{cnmc2026flex}, combined with a voluntary reduction of the railway permit to its verified need. The limiter, the battery and the control hierarchy are commissioned and tested against recorded and injected railway load profiles, including the loss of a neighbouring substation, without the railway depending on the outcome.

\emph{Stage 3: shared permit.} Under a regulatory sandbox or once the shared permit exists, the site is extended to its full size (\anDC\,MW with \anBatMW\,MW / \anBatMWh\,MWh in the model).

The critical path is administrative rather than technical: the access and connection permits, the environmental and municipal licences and the transmission bay typically take longer than the construction of the data centre. Land must lie outside the railway public domain and protection zones \cite{ley382015}, on railway-owned or private plots adjacent to the TSO substation.

\section{Security, resilience and train priority}
\label{sec:security}

\subsection{Design principle: no new failure mode for the railway}

The year 2026 has been the most difficult in the history of Spanish high-speed rail. On 18 January a derailment and subsequent collision at Adamuz, on the Madrid--Seville line, caused 46 deaths; preliminary findings point to a rail break at a weld and to deficiencies in inspection and alarm thresholds \cite{adamuz2026}. On 28 October, at Urda on the same line, a train ran into another that was stopped after a failure of the automatic signalling system had forced operation in a degraded manual mode; 69 people died, and the preliminary cause is a human error in traffic control during degraded operation \cite{urda2026}. The minister of transport and the president of the infrastructure manager left office in the following days, and a 50\,\% increase in railway investment over five years was announced \cite{infobae2026}. During the same year copper theft repeatedly disrupted high-speed traffic \cite{cable2026}, and in September an intrusion into the corporate systems of the infrastructure manager and the incumbent operator exfiltrated a large volume of customer data \cite{cyber2026}.

None of these events originated in the traction power supply. They are nonetheless the context in which any third-party use of railway electrical infrastructure will be judged, and they suggest three lessons that we adopt as requirements:
\begin{itemize}
\item[R1] \emph{Independence.} The data centre must not be able to cause, by failure, error or attack, any degradation of the railway supply. Whatever happens on the data-centre side, the railway must be exactly where it would be without it.
\item[R2] \emph{Fail-safe degraded modes.} The Urda sequence, in which the loss of an automatic system was followed by an unsafe manual fallback, is the pattern to exclude. On loss of measurement, communication or control, the data centre must reduce to a safe level automatically; no manual action may raise its limit.
\item[R3] \emph{No added burden.} The railway organisation, whose resources are being directed to safety and maintenance, must not acquire new operational tasks.
\end{itemize}
The following subsections show how the architecture of Section~\ref{sec:arch} meets them.

\subsection{Electrical priority of the train}

Requirement R1 is met electrically by construction: the two bays share the transmission busbar and nothing else. A fault in the data-centre installation is cleared by its own bay breaker, exactly as a fault in any other consumer connected to that substation; the railway bay, its protections and its transformers are untouched. The data centre cannot back-feed the catenary and does not use the traction transformers.

Priority in the use of the joint limit is guaranteed by the limiter relay, designed according to four rules:
\begin{enumerate}
\item It is a protection device, independent of the energy management system and of the data-centre IT network, with its own measurement of the data-centre bay and a hard-wired trip circuit.
\item Its setting is the joint limit agreed in the shared permit; it can only be changed by the TSO.
\item It is fail-safe: loss of the railway-power measurement, of auxiliary supply or of internal health trips the flexible block and holds the data centre at the firm level $D^{firm}$ of Eq.~\eqref{eq:firm}, which is safe for any railway load including N-1.
\item It is periodically tested by injection of recorded railway load profiles, and every operation is logged for the TSO and the regulator.
\end{enumerate}
With these rules the worst outcome of any malfunction is that the data centre loses load, never that the railway loses capacity. Requirement R3 follows: the railway does not operate, maintain or supervise anything new; the TSO already meters its bay.

\subsection{Contingencies and ride-through}

Table~\ref{tab:events} lists the events that affect the supply and the response of each party.

\begin{table}[t]
\caption{Supply events and response. Frequencies and durations are indicative.}
\label{tab:events}
\centering\footnotesize
\begin{tabular}{p{0.2\linewidth}p{0.13\linewidth}p{0.25\linewidth}p{0.3\linewidth}}
\toprule
Event & Duration & Effect on the data centre & Response\\
\midrule
Voltage dip from a transmission fault & 80--150\,ms & None if it rides through & UPS and grid-forming battery; no disconnection from the grid\\
Loss of a neighbouring traction substation & minutes to hours & Less headroom (Eq.~\ref{eq:n1}) & Battery, then curtailment of the flexible block\\
Loss of one data-centre transformer or feeder & hours & None & Second transformer takes the load\\
Loss of the transmission node & minutes to hours & Total loss of grid supply & Battery forms the 30\,kV island; generators start and supply the firm block\\
System-wide blackout & hours & Total loss of grid supply; trains are also stopped & As above; no conflict with the railway\\
\bottomrule
\end{tabular}
\end{table}

\emph{Loss of a neighbouring substation.} This is the contingency that sizes the curtailable block. With the batteries of the base case, the data centres sustain \mzbNoneSust\,MW of their \mzbBat\,MW on Madrid--Barcelona and \msNoneSust\,MW of \msBat\,MW on Madrid--Seville, so \mzbNoneCut\,MW and \msNoneCut\,MW must be shed (Fig.~\ref{fig:n1}); the median per site is \mzbNoneCutMed\,\% and \msNoneCutMed\,\%, within the 30\,\% assumed for the flexible block in all but three substations, where the site should be smaller or the flexible share larger. The curtailment can be staged over the minutes that the railway needs to reconfigure the catenary feeding, while the battery covers the transition.

\begin{figure}[t]
\centering
\includegraphics[width=\linewidth]{fig_n1}
\caption{Data-centre power sustained and curtailed at each site when a neighbouring traction substation is out of service (battery option). The marker shows the firm block.}
\label{fig:n1}
\end{figure}

\emph{Loss of grid supply.} The battery sized for train-priority shaving also bridges the start of the generators. At Anchuelo, \anBatMW\,MW / \anBatMWh\,MWh correspond to about \anBridgeMin\,min at rated power, against the tens of seconds that medium-voltage sets need to start and take load; reserving a minimum state of charge for this function is sufficient. In the blackout of 28 April 2025, caused by a combination of voltage and oscillation phenomena \cite{entsoe2026}, supply to the Iberian peninsula took many hours to restore and railway traffic stopped; in such an event there is no competition for capacity between the train and the data centre, and during restoration the limiter keeps the data centre at the firm level until the TSO releases the joint limit.

\subsection{Systemic risk for the transmission system}

Two dozen data centres along a line are a new type of risk for the TSO. In the Virginia event of 10 July 2024, about 1,500\,MW of data-centre load disconnected simultaneously because their protection and UPS logic reacted to a sequence of voltage dips \cite{nerc2025}; the sudden loss of load raised frequency and voltage. Loads of this kind sited on the same 400\,kV axis could behave in the same way, and the over-voltage sensitivity shown by the Iberian system in 2025 \cite{entsoe2026} makes the point more acute. The shared permit should therefore include grid-code obligations for the data centre, in line with the guidance being developed for large loads \cite{nerc2025}:
\begin{enumerate}
\item voltage ride-through: remain connected for dips within a defined envelope and count events on a rolling window instead of transferring to islanded supply after a fixed number of dips;
\item ramp-limited reconnection after any transfer to on-site supply;
\item ramp limits and filtering, by the battery, of workload-induced power oscillations \cite{li2024};
\item no use of generating sets to cover train-priority curtailment;
\item voltage support from the grid-forming battery \cite{lin2020}, which turns the site into an asset for the node.
\end{enumerate}

\subsection{Physical security}

Traction substations and their TSO counterparts are, by nature, isolated sites along the line. Co-location changes their physical risk profile in both directions. It adds permanent staffing, perimeter surveillance and lighting to places that are otherwise unattended, which is a deterrent against the copper theft that has repeatedly disrupted traffic \cite{cable2026}. It also adds a high-value target and more people and vehicles near railway assets. The design rules are those of separation: independent fenced perimeters and accesses for the railway substation, the TSO substation and the data centre; no shared cable routes, trenches or auxiliary supplies with railway installations; no data-centre construction or access through the railway public domain; and vetted personnel under the access rules applicable to critical-infrastructure operators \cite{ley82011,cer2022}. Fuel storage for the generators and the battery containers are sited and protected against fire with clearances to both substations.

\subsection{Cyber security}

The cyber risk has two faces. The first is that the data centre becomes a path into railway or TSO operational technology (OT). The second is that an attacker manipulates the data centre to overload the node. Both are closed by the same architecture, organised in zones and conduits \cite{iec62443}:
\begin{itemize}
\item \emph{One-way measurement.} The only information that crosses from the railway/TSO side is the active power of the railway bay. It is taken from the TSO metering at the bay and delivered through a unidirectional gateway (data diode) or as a hard-wired analogue signal. There is no routable connection between the data-centre networks and the railway control systems, signalling or telecommunications, and no connection at all to railway traffic management.
\item \emph{No command path.} The data centre receives no commands from the railway and sends none. Curtailment requests from the TSO use the channel already defined for flexible access permits.
\item \emph{Protection outside the attack surface.} The limiter is a hard-wired protection relay with local settings; compromise of the energy management system or of the IT workload scheduler can cause the data centre to consume \emph{less}, or to be tripped by the limiter, but not to exceed the joint limit.
\item \emph{Segmentation and standards.} The electrical OT of the data centre (substation automation under IEC~61850, battery and generator controls) is a separate zone from the IT infrastructure, with secured communications \cite{iec62351}, and both parties are subject to the obligations of essential entities \cite{nis2,ens2022}. On the railway side the applicable reference is CLC/TS~50701 \cite{ts50701}.
\end{itemize}
The September 2026 breach affected corporate IT and customer data, not operations \cite{cyber2026}; the literature has long warned that railway control systems \cite{lopez2015,kour2020} and power grids \cite{sun2018} are exposed when IT and OT are insufficiently separated. The proposal deliberately avoids creating any new IT/OT or OT/OT interface on the railway side.

\subsection{Risk register}

Table~\ref{tab:risk} summarises the threats considered, their potential impact on the railway and the mitigation. In every row the residual impact on the railway is nil or equal to the situation without a data centre, which is the design objective.

\begin{table}[t]
\caption{Risk register (design level).}
\label{tab:risk}
\centering\scriptsize
\begin{tabular}{p{0.23\linewidth}p{0.2\linewidth}p{0.33\linewidth}p{0.14\linewidth}}
\toprule
Threat & Potential impact & Mitigation & Residual impact on railway\\
\midrule
Data-centre load exceeds the headroom (control failure) & Overload of the joint limit & Independent limiter relay; battery; staged curtailment & None\\
Loss of railway-power measurement & Limit not enforced & Fail-safe: reduction to firm level & None\\
Electrical fault in the data-centre installation & Disturbance at the node & Own bay breaker and protections; selectivity study & As any other consumer at the node\\
Simultaneous trip of several data centres & Over-voltage, frequency rise & Ride-through and ramp obligations; grid-forming support & As for the system\\
Workload-induced power oscillations & Flicker, oscillation excitation & Battery filtering; ramp limits & None\\
Intrusion into data-centre IT or EMS & Malicious load changes & Hard-wired limiter; OT/IT segmentation & None\\
Pivot from data centre to railway OT & Compromise of railway systems & One-way gateway; no routable link; no command path & None (no interface)\\
Physical intrusion or theft & Damage to shared surroundings & Separate perimeters and routes; surveillance; vetted staff & Reduced (added presence)\\
Fire in battery or fuel store & Damage to adjacent substations & Clearances; fire protection; siting & Low\\
Traffic growth beyond forecast & Less headroom & Train has priority; periodic review of the limit; battery upgrade & None\\
\bottomrule
\end{tabular}
\end{table}

\section{Regulatory pathway, discussion and limitations}
\label{sec:discussion}

\subsection{What is possible today and what requires reform}
\label{sec:regulation}

Under the Spanish Electricity Act \cite{ley242013} and the access regulations \cite{rd11832020,cnmc2024circular}, an access permit belongs to one holder for its own consumption. The infrastructure manager cannot transfer or sublet capacity to another consumer, nor resell energy behind its meter; direct lines and closed distribution networks do not fit the case. Table~\ref{tab:reg} classifies the possible arrangements.

\begin{table}[t]
\caption{Regulatory routes for a data centre at a traction-substation node.}
\label{tab:reg}
\centering\footnotesize
\begin{tabular}{p{0.05\linewidth}p{0.42\linewidth}p{0.14\linewidth}p{0.27\linewidth}}
\toprule
 & Arrangement & Possible today & Capacity in the model\\
\midrule
A & Data centre behind the railway meter & No & --\\
B & Own permit and bay, flexible access (Type~3) & Yes & Node margin only\\
C & As B, plus voluntary reduction of the railway permit & Yes, with allocation risk & Firm: \mzbFirm{} / \msFirm\,MW\\
D & Data centre sharing the railway bay & No (dedicated bay required until 2029) & --\\
E & Shared access permit with priority & Requires reform & Up to \mzbBat{} / \msBat\,MW\\
\bottomrule
\end{tabular}
\end{table}

Route C is available now: the flexible access permits in force since September 2026 allow a consumer above 1\,MW to connect at transmission level accepting curtailment ordered by the TSO, with a dedicated bay \cite{cnmc2026flex}, and the 2026 decree-law opened a window to reduce granted capacity without forfeiting guarantees \cite{rdl72026}. Its weakness is that the capacity released by the railway returns to the node and is allocated by the general priority order or by tender; it cannot be directed to the co-located consumer. It also unlocks only the firm margin.

Route E is the proposal of this paper. It requires four changes: (i) a \emph{shared demand access permit} for several holders at the same node, with a joint limit and an agreed priority order, analogous to the hybridisation of generation facilities under one permit; (ii) acceptance that the curtailment signal of a flexible permit may be the load of the priority holder and not only an instruction from the system operator; (iii) allocation of capacity reduced by one consumer to a co-located consumer without tender, provided the node is not worse off; and (iv) explicit recognition that offering access headroom and land to non-railway users is compatible with the functions of the infrastructure manager, in line with the earlier recommendation of the competition authority on substations and renewable generation \cite{cnmc2023adif}. The definition of the joint limit should specify the averaging window, which Section~\ref{sec:results} shows to be decisive.

\subsection{Interaction with the data-centre decree}

The decree mandated by the 2026 decree-law had not been approved when this paper was completed (16 November 2026). The draft submitted to public hearing in August 2026 \cite{prd2026} applied to data centres of 1\,MW or more and required a PUE of 1.15, 80\,\% of consumption from new renewable plants and 80\,\% hourly matching. The revised draft sent to the Council of State, which issued a favourable opinion on 6 November \cite{mundiario2026}, raises the threshold to 10\,MW, relaxes the transitional PUE to 1.25, admits storage in the hourly matching and, most relevant here, introduces a transitional regime with lower matching targets for data centres that join the automatic power-reduction scheme of the system operator \cite{ashurst2026}. The regulator had asked for flexibility to be recognised \cite{cnmc2026ipn}. If confirmed, the decree will reward exactly the capability that the shared permit requires, and the result of Section~\ref{sec:results} quantifies it: flexibility reduces the renewable portfolio needed for compliance by about 30\,\%. The decree also affects the premise of scarcity: permits that are not built will lapse, which may free capacity at saturated nodes. The value of the proposal then rests less on avoiding a queue and more on reusing existing nodes, bays and land without new transmission assets.

\subsection{Why the two lines differ}

The busy line has little firm margin because its permits were sized to its peaks; its value lies in sharing in real time with storage. The quiet line has substantial firm margin today, because permits of similar size serve half the traffic; there the proposal could start under current rules. In both cases the capacity with batteries is of the order of 400\,MW per line, comparable to a large hyperscale campus but distributed over a dozen sites of 25--50\,MW, which suits the regional distribution of computing capacity and limits the systemic exposure of any single node.

\subsection{Limitations}

The results are preliminary in several respects, all stated as assumptions in Table~\ref{tab:assumptions}.

\emph{Load curves.} Substation loads are simulated, not measured. The real load curves of the infrastructure manager, at 1\,s resolution, are needed to validate peaks and ramps. The simulation uses approximate gradients, generic rolling-stock parameters, neutral sections at midpoints and no catenary losses; one day of one season was simulated.

\emph{Substation data.} Chainages are estimated ($\pm10$\,km), which shifts load between neighbours; six substations have no identified transmission bay and an assumed access; one granted value (12\,MW at Rueda de Jal\'on) is anomalously low. The fact that simulated peaks exceed 90\,\% of the granted access in \allTight{} substations must be reconciled with real data.

\emph{Contingencies.} The N-1 rule (half of the most loaded neighbour) is a simplification; the real redistribution depends on the feeding scheme and on the time needed to reconfigure the catenary, which also determines how fast the data centre must shed load.

\emph{Limit enforcement.} The resolution at which a joint limit would be monitored is a regulatory choice not yet made; the results span 1\,s to 15\,min.

\emph{Grid data.} Short-circuit powers are assumed; node margins for access beyond the current permits were not analysed.

\emph{Economics.} Unit costs and the sharing fee are order-of-magnitude assumptions; the renewable-matching analysis uses aggregate national profiles and is optimistic.

\emph{Timetable.} The model uses the October 2026 timetable. After the accident of 28 October services between Madrid and Andalusia were suspended and operating rules are under review; lower speeds or frequencies would reduce the traction load and increase headroom, so the Madrid--Seville results are conservative in that respect.

\emph{Security.} Section~\ref{sec:security} is a design-level analysis; a formal hazard analysis under railway safety standards and a site-specific security plan are required before implementation.

\subsection{Future work}

The immediate next step is validation with measured load curves and the real electrical data of each node, followed by site-level studies of individual substations (protection selectivity, power quality, short-circuit contribution). On the modelling side: a stochastic treatment of delays and traffic growth; co-optimisation of the battery for train-priority shaving, renewable matching and ancillary services; coordinated shifting of workloads between the sites of a line; and extension to the remaining high-speed lines. On the regulatory side: a legal design of the shared permit and of its metering and settlement.

\section{Conclusions}
\label{sec:conclusions}

High-speed traction substations in Spain hold transmission-level access that they use at \msUsePct--\mzbUsePct\,\% on average but fill for a few minutes at a time. With the access capacity actually granted to each railway bay and a timetable-driven simulation of two lines, this paper has shown that:

\begin{enumerate}
\item Under current rules, the capacity that the railway could release without losing N-1 firmness is small on the busy line (\mzbFirm\,MW on Madrid--Barcelona) and significant on the quiet one (\msFirm\,MW on Madrid--Seville).
\item A shared access permit with absolute train priority would allow \mzbFlex\,MW and \msFlex\,MW of data centres with a 30\,\% curtailable share, and \mzbBat\,MW and \msBat\,MW with power batteries of \mzbBatMWh\,MWh and \msBatMWh\,MWh, without any new transmission asset.
\item The averaging window of the joint limit is the most influential design parameter: 15-min instead of 1-min enforcement more than triples the storage-free capacity on the busy line.
\item The installation can be built with standard equipment: a dedicated bay at the same node, redundant transformers, a grid-forming battery limited to 2C, a two-block data centre and an independent fail-safe limiter. Its electrical cost is about 10\,\% of the data-centre investment.
\item The architecture gives the railway no new failure mode: the two installations share only a power limit, the only signal is a one-way measurement, every degraded mode reduces the data centre and never the train, and the railway acquires no operational duty. After the events of 2026 this is the condition for the proposal to be acceptable.
\item The flexibility needed to yield power to the train also lowers the cost of hourly renewable matching by about \resFlexValue\,\texteuro/MWh, and the latest draft of the Spanish data-centre decree moves towards rewarding it.
\end{enumerate}

The capacity is blocked by the design of access rules, not by the physics of the grid. Validating the simulated loads with measured data and testing the limiter at one pilot site are the two steps that would turn this estimate into an engineering project.


\section*{CRediT authorship contribution statement}
To be completed at submission following the CRediT taxonomy (conceptualisation, methodology, software, validation, writing, supervision).

\section*{Declaration of competing interest}
The authors declare that they have no known competing financial interests or personal relationships that could have appeared to influence the work reported in this paper. Any professional affiliation with railway or data-centre companies will be disclosed at submission.

\section*{Data availability}
The model, the input data that can be redistributed, the result tables and the scripts that generate every figure and every number quoted in this paper are available in the project repository (link to be provided at submission). The Renfe timetable feed is open data and must be downloaded separately as described in the repository.

\section*{Declaration of generative AI and AI-assisted technologies in the writing process}
During the preparation of this work the authors used a large language model assistant to support code development, literature search and drafting of the manuscript. After using this tool, the authors reviewed and edited the content as needed and take full responsibility for the content of the published article.

\section*{Funding}
This research did not receive any specific grant from funding agencies in the public, commercial, or not-for-profit sectors.

\bibliographystyle{elsarticle-num}
\bibliography{refs}

\end{document}
